\documentclass[11pt]{article}
\usepackage{geometry}
\geometry{a4paper, margin=1in}
\usepackage{amsmath,amssymb}
\usepackage{booktabs}
\usepackage{graphicx}
\usepackage{microtype}
\usepackage{url}
\usepackage[hidelinks]{hyperref}
\usepackage{xcolor}
\usepackage[htt]{hyphenat}
\sloppy

\title{Post-Tuning Beats Its Own Baseline and Still Loses to a Tuned Rule:\\
LLM Agents from 0.6B to 8B on a Simulated Beer Game}
\author{Pepijn Frenken\\
Independent researcher\\
\texttt{pepijnfrenken@mailbox.org}}
\date{2026-11-25}

\begin{document}
\maketitle

\begin{abstract}
A published claim reports that RL post-training turns a 4B language model into a cheaper supply-chain
agent than the same model untrained \cite{long2026}. We test a claim of that type at four sizes
(0.6B, 1.7B, 4B, 8B) under a pre-registered paired cost criterion, with matched untrained controls.
Post-training does help at 4B: 9{,}346 against 21{,}578, 30/30 paired wins. It does not help at
0.6B, where six trained rows lose and a seventh procedure stops itself; it makes 1.7B worse on every
seed (27{,}706 to 94{,}887 against 11{,}508); and at 8B one arm of four passes (14{,}594 against
89{,}633). The result that survives all of it is simpler: no arm at any size
beats the reference we actually care about. A self-tuning order-up-to rule costs 5{,}346 on the same
30 demand paths. The best arm is 1.75$\times$ that, the worst 17.8$\times$, and the ladder never
crosses it. Post-training moves the model within a band far above a 1970s inventory rule; scale did
not close the gap.
\end{abstract}

\section{What was tested}

\cite{long2026} reports a controlled beer-game experiment in which a GRPO-tuned 4B agent beats its
own untrained baseline on cost (v4 of \texttt{arXiv:2605.17036}, 21 September 2026, reports a mean
total cost falling from 1{,}585 to 952). We did not take the recipe on trust and we did not attempt a
verbatim re-run: we implemented an equivalent decision task, froze the evaluation rules before
training, and asked two questions that the published comparison leaves open.

\begin{enumerate}
  \item \textbf{Does the effect survive at other sizes?} The claim is about one model. We ran the same
  task at 0.6B, 1.7B, 4B and 8B.
  \item \textbf{Does the tuned agent beat the algorithm?} Beating an untrained copy of itself is a low
  bar. The reference that matters in inventory control is a tuned order-up-to rule, and our
  environment computes one per demand path.
\end{enumerate}

The short answer to the first question is: at one size, and fragile. The short answer to the second is
no, at every size, by a wide margin. Section~\ref{sec:results} gives the numbers and
Section~\ref{sec:solid} states exactly what they do and do not establish.

\section{Environment and the two baselines}\label{sec:env}

The simulator is the four-echelon beer game used throughout this project: four roles (retailer,
wholesaler, distributor, factory), a two-week shipping delay, holding and backlog costs of 1 and 2 per
unit per week, a 36-week horizon, zero initial inventories, and stochastic demand. The agent sees the
local state and returns one order quantity per decision. The action set is a fixed grid of 26 values
(0--30, then 36--200), read as 26 label tokens, so a decision is a single-letter choice rather than
free text.

Two baselines are in play, and keeping them apart is the point of this paper.

\textbf{The untrained self.} The same checkpoint, same prompt, same grid, no training. This is the
baseline the published claim uses, and it is a fair test of whether training did anything.

\textbf{The tuned rule.} A walk-forward order-up-to policy that re-estimates its parameters from the
first 18 weeks of each demand path and then plays the remaining horizon. It is computed inside the same
simulator with the same cost function, and it appears in three versions on the same 30 noisy demand
paths: tuned ($\theta,\lambda$ fitted), untuned (fixed parameters), and an oracle that knows the
demand. Medians: tuned 5{,}346, untuned 5{,}423, oracle 4{,}594. Tuning therefore buys almost nothing
on noisy demand; the rule's value is in the rule, not in the fitting.

Comparing a model to this rule is the comparison a practitioner would make, and it is the one the
published claim does not make.

\section{Protocol}\label{sec:protocol}

Every evaluation rule was committed before training and is unchanged since. The full definitions,
thresholds per size and per-round read schedule are in the methods supplement \cite{supplement}
(\S1--\S2); what follows is the part needed to read the results.

\begin{itemize}
  \item \textbf{The criterion (G1).} On 30 shared noisy demand paths (seeds 1000--1029), the trained
  arm must reduce cost by at least 15\% against its matched untrained control on two statistics (ratio
  of medians, and median of per-path proportional reductions), win at least 20 of the 30 paired paths,
  and clear a one-sided sign test at $p<0.02$. A size passes if at least two of its three seeds pass.
  \item \textbf{The stop-loss.} After each training round the policy's output distribution is
  re-measured. A collapse branch (median entropy below the frozen low bound with modal mass above
  threshold, or entropy above the high bound with an information statistic below threshold) fires on
  two consecutive reads and stops the arm. The record states the arm stopped; it does not quietly
  finish the schedule.
  \item \textbf{The pooling rule.} Arms are selected by a pooled mean over the seeds run at that
  learning rate, fixed before outcomes were reviewed. It is why one favourable seed at 8B does not
  become the headline.
  \item \textbf{The ledger.} 24 planned tasks (six wave-7 comparators at 0.6B, eighteen wave-8 R1
  tasks), each with a terminal disposition read from its own markers. 12 carry a matched paired final
  evaluation; 12 do not (6 stopped, 4 trained without one, 2 unattempted).
\end{itemize}

Training is LoRA ($r{=}8$, $\alpha{=}16$, attention projections), batch 32, 170 optimizer steps per
round, three rounds, exact-enumeration policy gradient against a frozen reference KL. Section~A lists
the deltas from the published recipe.

\section{Results}\label{sec:results}

Table~\ref{tab:rungs} is the whole ladder. Costs are medians over the same 30 noisy demand paths unless
marked otherwise; ``wins'' counts paths on which the trained arm is cheaper than its matched untrained
control.

\begin{table}[ht]
\centering
\caption{Four sizes against their own untrained baseline, matched paired criterion. October and
November cohorts are separate entries because they are separate evidence classes; see
Section~\ref{sec:solid}.}
\label{tab:rungs}
\small
\begin{tabular}{@{}llrrr@{}}
\toprule
Size & Cohort & Untrained median & Trained median & Paired wins \\
\midrule
0.6B & wave-7 SFT, RLCD (6 arm--seed rows) & 11{,}569 & 20{,}594--37{,}384 & 0/30 each \\
1.7B & November & 11{,}508 & 27{,}706 / 79{,}094 / 94{,}887 & 7/30, 0/30, 0/30 \\
4B & November & 21{,}578 & 9{,}346 / 17{,}958 & 30/30, 25/30 \\
8B & October survivor & 89{,}633 & 14{,}594 & 29/30 \\
\midrule
1.7B & October (recorded) & 11{,}508 & 21{,}468 / 43{,}660 / 71{,}270 & reported 0/30 \\
4B & October (recorded) & 21{,}578 & 8{,}315 / 7{,}612 & reported 30/30 \\
\bottomrule
\end{tabular}
\end{table}

\textbf{4B, the size where training works.} The November rebuild passes G1 on both seeds: 9{,}346
(30/30 wins, $p = 9.3\times10^{-10}$) and 17{,}958 (25/30, $p \approx 1.6\times10^{-4}$). The
reduction is real and reproducible from committed per-run rows. It is also seed-dependent: the two
seeds differ by $1.52\times$ in median and 1.42$\times$ in mean, so the arm-level mean
(13{,}152) is what we report, not the flattering seed.

\textbf{1.7B, where training is worse than doing nothing.} All three November seeds lose to their
untrained control: 27{,}706 (7/30 wins), 79{,}094 and 94{,}887 (0/30). The effect is not merely absent;
it is adverse by 2.4$\times$ to 8.2$\times$. Training at this size found a policy that is confidently
wrong, and the collapse detector watched it happen.

\textbf{0.6B, where the procedure stops itself.} Wave-7 produced the closed-choice readout that makes
the whole study measurable, and six trained arm--seed rows that all lose (0/30 wins each). Wave-8's R1
arms tripped the collapse branch on all six tasks; five carry persisted stop markers and three tasks
were resumed past a stop, two of them through one unattended resubmission.

\textbf{8B, where one arm passes and three do not.} The $10^{-5}$ seed-1 checkpoint survives and passes
G1 against a matched untrained-8B control: 14{,}594 against 89{,}633, 29/30, $p = 2.9\times10^{-8}$.
It is not the arm the frozen selection rule chose. The selected $5\times10^{-5}$ rate was retrained in
November and reads $0.108$ and $0.164$ of a \emph{curriculum} control (10{,}097 and 15{,}355 against
93{,}603, 30/30 and 28/30), which is favourable but is not the registered comparison: those legs ran
the curriculum demand pattern and were paired against a noisy control, so the two sides share no demand
path. An earlier draft of this paper claimed a G1 pass for them. That claim was withdrawn; the correct
statement is that the selected rate's paired outcome is unestablished, not that it failed.

\subsection{The comparison that matters}

Table~\ref{tab:formula} puts every arm beside the tuned rule on the same 30 demand paths.

\begin{table}[ht]
\centering
\caption{Every model arm against the tuned walk-forward rule, medians of total cost over the same 30
noisy demand paths. Ratio is arm median divided by rule median. The comparison is distributional, not
paired: see Section~\ref{sec:solid}.}
\label{tab:formula}
\small
\begin{tabular}{@{}lrr@{}}
\toprule
Arm & Median cost & Ratio to tuned rule (5{,}346) \\
\midrule
Tuned order-up-to rule & 5{,}346 & 1.00 \\
Oracle rule (knows demand) & 4{,}594 & 0.86 \\
\midrule
4B post-trained, seed 0 & 9{,}346 & 1.75 \\
8B post-trained, October survivor & 14{,}594 & 2.73 \\
4B untrained & 21{,}578 & 4.04 \\
1.7B post-trained, seed 0 & 27{,}706 & 5.18 \\
1.7B post-trained, seed 1 & 79{,}094 & 14.80 \\
1.7B post-trained, seed 2 & 94{,}887 & 17.75 \\
8B untrained & 89{,}633 & 16.77 \\
\bottomrule
\end{tabular}
\end{table}

Not one trained arm, at any size, comes within 75\% of the rule. The best is the 4B arm that passes
G1, and it is 1.75$\times$ the rule while beating its own untrained self by 2.3$\times$. Both
statements are true and they point in opposite directions: post-training is a real intervention on the
model, and the model is still nowhere near the algorithm. Scaling from 0.6B to 8B did not close the
gap; at 1.7B it widened it.

\section{What this does and does not establish}\label{sec:solid}

\textbf{Established.} The published comparison type reproduces at 4B: post-training roughly halves cost
against the same model untrained, on matched demand paths, with a sign test. It does not reproduce as a
general property of scale: one size of four improves, one is neutral, one is sharply adverse, and one
passes on a single unselected arm. And no arm approaches the tuned rule.

\textbf{Not established, and stated as such.} The October 4B and 1.7B numbers are reported because
earlier drafts published them, not because we can recompute them. The October verdict document records
two 4B costs (8{,}315 and 7{,}612) and a control value, with no win count and no per-path reduction; the
per-run files behind both October cohorts were never committed and did not survive. The November
rebuild supersedes them as evidence, and every verdict in this paper rests on November rows or on the
8B October adapter, which does.

\textbf{The formula comparison is distributional.} \texttt{results/wave7\_floors.json} gives the rule's
cost per demand path, and the arm files give cost per run, but we could not verify a common path
identifier between the two. As a check, we correlated arm cost with rule cost across indices and got
$-0.07$ to $+0.16$, which is noise; a real paired join should show positive correlation, since harder
demand should cost both more. We therefore compare medians over the same 30-path band and claim no
per-path dominance.

\textbf{One arm is not a size.} 8B is represented by a single checkpoint that the frozen selection rule
did not choose, whose adapters no longer exist, evaluated in November on a Blackwell card with a newer
stack than the October legs. Its pass is real; it is not a rung verdict.

\textbf{Variance is training-side, not evaluation-side.} Two independent 4B evaluation jobs on different
days produced byte-identical outputs, and the untrained band control returned the same value to five
decimals twice. The 14--39\% differences between October and November therefore live in training, which
we can point at but cannot yet explain.

\textbf{Other limits.} Simulated environment; one demand family for the headline comparisons; LoRA and
an exact-enumeration objective rather than the published GRPO recipe; effective batch matched across
hardware but optimizer-step and numerical equivalence not tested; 24 planned tasks is too few for any
scaling claim, and none is made.

\section{Provenance and what is missing}\label{sec:prov}

Every number here is labelled by cohort. November rows regenerate from committed per-run files:
paired cost rows, summaries, round readings and frozen gate files under \texttt{results/}. October
numbers are recorded, not recomputable, and the supplement's S5 lists each one with the file that does
and does not hold it. The submission tag is \texttt{arxiv-submission-1}; the protocol, the ledger, the
detector readings, the failures and the losses are in the supplement \cite{supplement}.

Two disclosures belong in the body rather than a footnote. First, an earlier draft of this paper
claimed a November G1 pass for the selected 8B rate; the pairing was wrong and the claim is withdrawn,
which is why Section~\ref{sec:results} reports it as band-level and unestablished. Second, a transfer of
evaluation traces through a notebook channel corrupted the files; they were deleted rather than patched,
and the 30-row cost extracts that replaced them are what the G1 reads consume.

\appendix
\section{Deltas from the published recipe}\label{app:deltas}

Adaptation, not replication. Differences that could plausibly matter: closed-choice readout over 26
label tokens instead of free-form text; exact-enumeration policy gradient with a frozen per-state KL
anchor instead of group-sampled REINFORCE; LoRA instead of full fine-tuning; a 36-week horizon with our
cost constants; and an explicit collapse stop-loss the published recipe does not describe. Shared with
the reference: the frozen reference-policy KL term, the simulator's structure, and the comparison
against an untrained copy of the same model.

\begin{thebibliography}{9}\itemsep 0pt
\bibitem{long2026} C.~X. Long, D.~Simchi-Levi, F.~Zhu, H.~Su, A.~P.~Calmon,
F.~P.~Calmon, ``Reliability and Effectiveness of Autonomous AI Agents in Supply Chain
Management,'' \texttt{arXiv:2605.17036}, version 4 (21 September 2026); accessed 25 November 2026.
\bibitem{supplement} P.~Frenken, ``Methods supplement: gates, per-round readings, task ledger and
provenance inventory for the four-scale closed-choice bullwhip study,'' 2026. Committed with this
paper at tag \texttt{arxiv-submission-1}.
\end{thebibliography}

\end{document}
